% Problema 3, parte 1: campi da incollare nella piattaforma (uno per sezione)

% ===== CAMPO: 1. Result and scope =====
Full written proof of all parts of cell 1: (i) for n=T\_k every
partition reaches δ\_k and δ\_k is the unique cyclic partition; (ii) for
n=T\_\{k-1\}+r the cyclic partitions are exactly the C(k,r) partitions
λ\^{}ε (positive numbers among k-i+ε\_i, ε∈\{0,1\}\^{}k of weight r), on
which B acts as cyclic rotation of ε; (iii) the number of cycles is
(1/k)Σ\_\{d\textbar gcd(k,r)\} φ(d) C(k/d, r/d). The convergence step is
proved via an ordered dynamics, a diagonal potential Φ with exact
equality characterisation, and a hole-chasing lemma. Exact brute-force
check agrees for all n≤40.


% ===== CAMPO: 2. Proof =====
\subparagraph{Notation}\label{notation}

Throughout \(n\ge 1\) has rank \(k\), i.e.~\(n=T_{k-1}+r\) with
\(1\le r\le k\) (equivalently \(T_{k-1}<n\le T_k\)). For a \(0/1\)-word
\(\varepsilon=(\varepsilon_1,\dots,\varepsilon_k)\) of weight
\(|\varepsilon|=\sum\varepsilon_i=r\) put
\[\lambda^{\varepsilon}:=\text{the partition whose parts are the positive numbers among } k-i+\varepsilon_i,\quad i=1,\dots,k,\]
\[W_{k,r}:=\{\varepsilon\in\{0,1\}^k:|\varepsilon|=r\},\qquad S_{k,r}:=\{\lambda^{\varepsilon}:\varepsilon\in W_{k,r}\}.\]

\textbf{Fact 0.} \(\lambda^\varepsilon\) is a partition of \(n\); the
sequence \((k-i+\varepsilon_i)_{i=1}^k\) is weakly decreasing, its first
\(k-1\) entries are \(\ge1\), and the \(k\)-th is \(\varepsilon_k\).
Moreover \(\varepsilon\mapsto\lambda^\varepsilon\) is injective on
\(W_{k,r}\), and \(S_{k,k}=\{\delta_k\}\).

\emph{Proof.} Sum: \(\sum_i(k-i)+\sum_i\varepsilon_i=T_{k-1}+r=n\).
Monotone:
\((k-i+\varepsilon_i)-(k-i-1+\varepsilon_{i+1})=1+\varepsilon_i-\varepsilon_{i+1}\ge0\).
For \(i\le k-1\), \(k-i+\varepsilon_i\ge1\). Injective:
\(\lambda^\varepsilon\) has \(k-1\) or \(k\) parts; pad it with zeros to
length exactly \(k\); the padded weakly decreasing sequence of a
multiset is unique, and \((k-i+\varepsilon_i)_i\) is such a sequence, so
the \(i\)-th padded part equals \(k-i+\varepsilon_i\) and
\(\varepsilon_i\) is recovered. For \(r=k\) the only word is
\((1,\dots,1)\) and \(\lambda^\varepsilon=(k,k-1,\dots,1)=\delta_k\).
\(\square\)

\subparagraph{\texorpdfstring{Part I. \(B\) acts on \(S_{k,r}\) as a
rotation}{Part I. B acts on S\_\{k,r\} as a rotation}}\label{part-i.-b-acts-on-s_kr-as-a-rotation}

\textbf{Lemma 1.} For \(\varepsilon\in W_{k,r}\),
\(B(\lambda^\varepsilon)=\lambda^{\sigma\varepsilon}\) where
\(\sigma\varepsilon:=(\varepsilon_k,\varepsilon_1,\varepsilon_2,\dots,\varepsilon_{k-1})\).
In particular \(S_{k,r}\) is \(B\)-invariant, \(B^k\) is the identity on
\(S_{k,r}\), and every element of \(S_{k,r}\) is cyclic.

\emph{Proof.} By Fact 0, \(\lambda^\varepsilon\) has
\(s=k-1+\varepsilon_k\) parts: the parts \(k-i+\varepsilon_i\) for
\(i\le k-1\), plus the part \(1\) (coming from \(i=k\)) iff
\(\varepsilon_k=1\). Applying \(B\): each part \(k-i+\varepsilon_i\)
(\(i\le k-1\)) becomes \(k-(i+1)+\varepsilon_i\), which is positive for
\(i\le k-2\) and equals \(\varepsilon_{k-1}\) for \(i=k-1\) (kept iff it
is \(1\)); the part \(1\) present when \(\varepsilon_k=1\) becomes \(0\)
and is discarded; the new part is \(s=k-1+\varepsilon_k\). Hence
\(B(\lambda^\varepsilon)\) has as parts the positive numbers among
\[k-1+\varepsilon_k\ \ (\text{index }i'=1),\qquad k-i'+\varepsilon_{i'-1}\ \ (i'=2,\dots,k),\]
which are exactly the numbers \(k-i'+(\sigma\varepsilon)_{i'}\),
\(i'=1,\dots,k\). So
\(B(\lambda^\varepsilon)=\lambda^{\sigma\varepsilon}\). Since \(\sigma\)
preserves weight and \(\sigma^k=\mathrm{id}\),
\(B^k(\lambda^\varepsilon)=\lambda^{\sigma^k\varepsilon}=\lambda^\varepsilon\)
with \(k\ge1\), so \(\lambda^\varepsilon\) is cyclic. \(\square\)

\subparagraph{\texorpdfstring{Part II. Every partition of \(n\) enters
\(S_{k,r}\)}{Part II. Every partition of n enters S\_\{k,r\}}}\label{part-ii.-every-partition-of-n-enters-s_kr}

II.1 Ordered dynamics

A \emph{composition} is a finite sequence \(c=(c_1,\dots,c_s)\) of
positive integers. Define
\[\tilde B(c):=\text{the sequence }(s,\ c_1-1,\ \dots,\ c_s-1)\text{ with all zero entries deleted}.\]
The multiset of entries of \(\tilde B(c)\) is
\(\{s\}\cup\{c_i-1:c_i\ge2\}\), which is by definition the multiset of
parts of \(B(\mathrm{sort}(c))\). Hence, starting from
\(c^{(0)}:=\lambda\) (as a sequence) and putting
\(c^{(t+1)}:=\tilde B(c^{(t)})\), we have
\[B^t(\lambda)=\mathrm{sort}(c^{(t)})\quad\text{for all }t\ge0. \tag{1}\]
Write \(s_t\) for the length of \(c^{(t)}\). The \emph{cell set} of a
composition is
\(C(c):=\{(i,j):1\le i\le s,\ 1\le j\le c_i\}\subset\mathbb Z_{\ge1}^2\);
put \(C_t:=C(c^{(t)})\). Cell \((i,j)\) lies on \emph{diagonal} \(i+j\).
Two properties hold for every composition, by definition:
(\textbf{closed}) if \((i,j)\in C(c)\) and \(1\le j'\le j\) then
\((i,j')\in C(c)\); (\textbf{gap-free}) if \((i,1)\in C(c)\) and
\(1\le i'\le i\) then \((i',1)\in C(c)\); and \((i,1)\notin C(c)\) iff
\(i>s\).

II.2 The potential

For a finite sequence \(e=(e_1,\dots,e_m)\) of non-negative integers put
\(\Phi(e):=\sum_{i=1}^m\sum_{j=1}^{e_i}(i+j)\in\mathbb Z_{\ge0}\) (zero
entries contribute nothing).

\textbf{Lemma 2.} For every composition \(c\),
\(\Phi(\tilde B(c))\le\Phi(c)\), with equality if and only if the set
\(\{i: c_i=1\}\) is a terminal segment of \(\{1,\dots,s\}\)
(i.e.~\(c_i=1\Rightarrow c_j=1\) for all \(j>i\)).

\emph{Proof.} Let \(e:=(s,c_1-1,\dots,c_s-1)\) (length \(s+1\), zeros
allowed). Then
\[\Phi(e)=\sum_{j=1}^{s}(1+j)+\sum_{i=1}^{s}\sum_{j=1}^{c_i-1}(i+1+j)=\sum_{i=1}^{s}(i+1)+\sum_{i=1}^{s}\sum_{j'=2}^{c_i}(i+j')=\sum_{i=1}^{s}\sum_{j'=1}^{c_i}(i+j')=\Phi(c),\]
where we substituted \(j'=j+1\) and used that the term \(j'=1\) of
column \(i\) is \(i+1\). Now \(\tilde B(c)\) is obtained from \(e\) by
deleting its zero entries; an entry \(e_q>0\) preceded by \(z_q\) zeros
moves to position \(q-z_q\), so \(\Phi\) decreases by
\(\sum_{q:e_q>0}z_q e_q\ge0\), with equality iff no positive entry is
preceded by a zero, i.e.~iff the zeros of \(e\) form a terminal segment.
Since \(e_1=s\ge1\) and \(e_{i+1}=0\iff c_i=1\), this says exactly that
\(\{i:c_i=1\}\) is terminal in \(\{1,\dots,s\}\). \(\square\)

By Lemma 2, \((\Phi(c^{(t)}))_{t\ge0}\) is a non-increasing sequence of
non-negative integers, hence eventually constant: there is \(t_0\) with
\(\Phi(c^{(t+1)})=\Phi(c^{(t)})\) for all \(t\ge t_0\). By the equality
case of Lemma 2:
\[\text{for all }t\ge t_0:\quad \{i:c^{(t)}_i=1\}\text{ is a terminal segment of }\{1,\dots,s_t\}. \tag{2}\]

II.3 Diagonal motion in the steady regime

Define \(\rho:\mathbb Z_{\ge1}^2\to\mathbb Z_{\ge1}^2\) by
\(\rho(i,j)=(i+1,j-1)\) if \(j\ge2\) and \(\rho(i,1)=(1,i)\). Then
\(\rho\) preserves \(i+j\), and \(\rho\) is a bijection (inverse:
\((i,j)\mapsto(i-1,j+1)\) for \(i\ge2\), \((1,j)\mapsto(j,1)\)). On
diagonal \(d\ge2\), whose cells are \((i,d-i)\), \(1\le i\le d-1\),
\(\rho\) sends column \(i\) to column \(i+1\) for \(i\le d-2\) and
column \(d-1\) to column \(1\); hence \(\rho^m\) sends the cell of
diagonal \(d\) in column \(i\) to the cell in the unique column
\(\equiv i+m \pmod{d-1}\) in \(\{1,\dots,d-1\}\).

\textbf{Lemma 3.} For all \(t\ge t_0\): \(C_{t+1}=\rho(C_t)\).

\emph{Proof.} Let \(c=c^{(t)}\), \(e=(s,c_1-1,\dots,c_s-1)\) as in Lemma
2. The cells of \(e\) (defined as for compositions, zero entries giving
no cells) are \(\{(1,j):1\le j\le s\}=\rho(\{(j,1):1\le j\le s\})\)
together with
\(\{(i+1,j-1):2\le j\le c_i\}=\rho(\{(i,j)\in C_t: j\ge2\})\); so the
cell set of \(e\) is \(\rho(C_t)\). By (2), the zeros of \(e\) form a
terminal segment, so deleting them does not move any positive entry, and
\(C_{t+1}=C(\tilde B(c))\) is the cell set of \(e\). \(\square\)

Consequently, for \(t\ge t_0\) and \(m\ge0\), \(C_{t+m}=\rho^m(C_t)\),
and since \(\rho^m\) is injective,
\[x\in C_t\iff\rho^m(x)\in C_{t+m}\qquad(x\in\mathbb Z_{\ge1}^2). \tag{3}\]

II.4 The hole lemma

\textbf{Lemma 4.} Let \(t\ge t_0\) and \(d\ge2\). If some cell of
diagonal \(d\) is not in \(C_t\) (a \emph{hole}), then no cell of any
diagonal \(D>d\) lies in \(C_t\).

\emph{Proof.} Diagonal \(2\) consists of the single cell
\((1,1)\in C_t\) (as \(n\ge1\), \(s_t\ge1\)), so \(d\ge3\). Suppose
\((p,d-p)\notin C_t\) and \((q,D-q)\in C_t\) with \(D>d\). For \(m\ge0\)
let \(h_m\in\{1,\dots,d-1\}\) be the column of \(\rho^m(p,d-p)\) and
\(\kappa_m\in\{1,\dots,D-1\}\) the column of \(\rho^m(q,D-q)\); by (3),
\(\rho^m(p,d-p)\notin C_{t+m}\) and \(\rho^m(q,D-q)\in C_{t+m}\). By
II.3, \(h_{m+1}\equiv h_m+1\pmod{d-1}\) and
\(\kappa_{m+1}\equiv\kappa_m+1\pmod{D-1}\).

\emph{Constraint (ii).} If \(h_m=d-1\) then \(\kappa_m\le d-2\). Indeed
the hole is then the cell \((d-1,1)\notin C_{t+m}\), so \(d-1>s_{t+m}\)
(II.1), i.e.~\(s_{t+m}\le d-2\), and every cell of \(C_{t+m}\) has
column \(\le d-2\).

Since \(h_m\) runs cyclically through \(1,\dots,d-1\), there are
infinitely many \(m\) with \(h_m=d-1\). Take one such \(m\) and put
\(u:=\kappa_m\in\{1,\dots,d-2\}\) (by (ii)). For \(l=1,\dots,d-1\) we
have \(h_{m+l}=l\) (the hole goes from column \(d-1\) to
\(1,2,\dots,d-1\)), and \(\kappa_{m+l}=u+l\) as long as \(u+l\le D-1\).

\emph{Case A: \(u\le D-d\).} Then \(u+(d-1)\le D-1\), so
\(\kappa_{m+d-1}=u+d-1\ge d\), while \(h_{m+d-1}=d-1\): this contradicts
(ii).

\emph{Case B: \(u\ge D-d+1\).} Then \(D-u\le d-1\); the card is in
column \(D-1\) at \(l=D-1-u\), in column \(1\) at \(l=D-u\), and in
column \(l-(D-u)+1\) for \(D-u\le l\le d-1\) (no further wrap, since
\(l-(D-u)+1\le d-D+u\le u\le d-2<D-1\)). At \(l=d-1\) the hole is again
in column \(d-1\) and the card in column \(u':=u-(D-d)\), with
\(1\le u'\le u-1\) because \(D-d\ge1\) and \(u\ge D-d+1\).

So each time the hole is in column \(d-1\), either Case A gives a
contradiction, or Case B occurs and at the next such time the card's
column is a strictly smaller positive integer. Case B cannot occur
forever (positive integers cannot decrease strictly infinitely often),
so Case A occurs at some time --- a contradiction. Hence no such card
exists. \(\square\)

II.5 Conclusion of Part II

\textbf{Theorem A.} For every partition \(\lambda\) of \(n\) (rank
\(k\), \(n=T_{k-1}+r\), \(1\le r\le k\)) there is \(t\) with
\(B^t(\lambda)\in S_{k,r}\). In particular for \(n=T_k\) there is \(t\)
with \(B^t(\lambda)=\delta_k\).

\emph{Proof.} Take \(t=t_0\) and \(C:=C_{t}\). Diagonal \(2\) is full
(\(=\{(1,1)\}\subseteq C\)); \(C\) is finite, so
\(K:=\max\{m\ge2:\text{diagonals }2,\dots,m\text{ are all contained in }C\}\)
exists, and diagonal \(K+1\) has a hole. By Lemma 4 no cell of diagonal
\(\ge K+2\) is in \(C\). Therefore
\[C=\{(i,j):i+j\le K\}\ \cup\ R,\qquad R\subseteq\{(i,K+1-i):1\le i\le K\},\quad r':=|R|\le K-1.\]
Set \(\varepsilon_i:=1\) if \((i,K+1-i)\in R\) and \(0\) otherwise
(\(1\le i\le K\)). Column \(i\le K\) of \(C\) consists of
\((i,1),\dots,(i,K-i)\) plus possibly \((i,K+1-i)\), so it has height
\(K-i+\varepsilon_i\); columns \(>K\) are empty. Hence \(c^{(t)}\) is
the sequence of positive numbers among \(K-i+\varepsilon_i\)
(\(i=1,\dots,K\)), and \(n=|C|=T_{K-1}+r'\) with \(0\le r'\le K-1\). If
\(r'\ge1\): then \(T_{K-1}<n\le T_K\), so \(K=k\), \(r'=r\), and by (1)
\(B^t(\lambda)=\mathrm{sort}(c^{(t)})=\lambda^\varepsilon\in S_{k,r}\)
(the sequence is already weakly decreasing by Fact 0). If \(r'=0\): then
\(n=T_{K-1}\), so \(k=K-1\), \(r=k\), and
\(c^{(t)}=(K-1,K-2,\dots,1)=\delta_k\),
i.e.~\(B^t(\lambda)=\delta_k\in S_{k,k}\). \(\square\)

\subparagraph{Part III. Cyclic partitions and
cycles}\label{part-iii.-cyclic-partitions-and-cycles}

\textbf{Theorem B.} The cyclic partitions of \(n=T_{k-1}+r\)
(\(1\le r\le k\)) are exactly the elements of \(S_{k,r}\), i.e.~the
partitions \(\lambda^\varepsilon\) with \(\varepsilon\in\{0,1\}^k\) of
weight \(r\) --- explicitly, the partitions whose parts are the positive
numbers among
\(k-1+\varepsilon_1,\ k-2+\varepsilon_2,\ \dots,\ 1+\varepsilon_{k-1},\ \varepsilon_k\)
with exactly \(r\) of the \(\varepsilon_i\) equal to \(1\). There are
\(\binom kr\) of them. For \(n=T_k\) (\(r=k\)) the only cyclic partition
is \(\delta_k\), and every partition of \(T_k\) reaches \(\delta_k\).

\emph{Proof.} (\(\supseteq\)) Lemma 1. (\(\subseteq\)) Let \(\lambda\)
be cyclic, \(B^i(\lambda)=\lambda\) with \(i\ge1\). By Theorem A,
\(B^t(\lambda)\in S_{k,r}\) for some \(t\). Choose an integer \(q\) with
\(qi\ge t\). Then
\(\lambda=B^{qi}(\lambda)=B^{qi-t}\big(B^t(\lambda)\big)\in S_{k,r}\) by
\(B\)-invariance of \(S_{k,r}\) (Lemma 1). The count \(\binom kr\) is
Fact 0 (injectivity) and \(|W_{k,r}|=\binom kr\). The triangular
statements are the case \(r=k\) (\(S_{k,k}=\{\delta_k\}\)) together with
Theorem A. \(\square\)

\textbf{Theorem C.} The number of distinct cycles of \(B\) on the
partitions of \(n=T_{k-1}+r\) is
\[N(k,r)=\frac1k\sum_{d\mid\gcd(k,r)}\varphi(d)\binom{k/d}{r/d},\] where
\(\varphi\) is Euler's function. For \(n=T_k\) this equals \(1\).

\emph{Proof.} Every cycle of \(B\) consists of cyclic partitions, hence
(Theorem B) lies in \(S_{k,r}\), and every element of \(S_{k,r}\) lies
on a cycle (Lemma 1). By Lemma 1 and the injectivity of
\(\varepsilon\mapsto\lambda^\varepsilon\), the cycles of \(B\) on
\(S_{k,r}\) correspond bijectively to the orbits of the cyclic group
\(G=\langle\sigma\rangle=\{\sigma^j:0\le j\le k-1\}\) (of order \(k\))
acting on \(W_{k,r}\) by rotation.

\emph{Orbit count (Burnside, with proof).} Count pairs
\((g,\varepsilon)\in G\times W_{k,r}\) with \(g\varepsilon=\varepsilon\)
in two ways:
\(\sum_{g\in G}|\mathrm{Fix}(g)|=\sum_{\varepsilon}|\mathrm{Stab}(\varepsilon)|=\sum_{\varepsilon}\frac{|G|}{|G\varepsilon|}=|G|\cdot\#\text{orbits}\)
(orbit--stabiliser; each orbit \(O\) contributes
\(\sum_{\varepsilon\in O}|G|/|O|=|G|\)). Hence
\(\#\text{orbits}=\frac1k\sum_{j=0}^{k-1}|\mathrm{Fix}(\sigma^j)|\).

\emph{Fixed points of \(\sigma^j\).} \(\sigma^j\) shifts positions by
\(j\) modulo \(k\); its cycles on \(\{1,\dots,k\}\cong\mathbb Z_k\) are
the cosets \(i+g\mathbb Z_k\) where \(g=\gcd(j,k)\) (because
\(\{jm\bmod k\}=g\mathbb Z_k\)), so there are \(g\) cycles each of
length \(k/g\). A word is fixed by \(\sigma^j\) iff it is constant on
each cycle; choosing \(a\) cycles to carry \(1\)'s gives weight
\(a\,k/g\). So \(|\mathrm{Fix}(\sigma^j)|=\binom{g}{rg/k}\) if
\((k/g)\mid r\), and \(0\) otherwise. The number of
\(j\in\{0,\dots,k-1\}\) with \(\gcd(j,k)=g\) is \(\varphi(k/g)\) (write
\(j=gj'\), \(0\le j'<k/g\), \(\gcd(j',k/g)=1\)). Putting \(d=k/g\) (so
\(d\mid k\), and the term is nonzero iff \(d\mid r\)):
\[N(k,r)=\frac1k\sum_{d\mid k,\ d\mid r}\varphi(d)\binom{k/d}{r/d}=\frac1k\sum_{d\mid\gcd(k,r)}\varphi(d)\binom{k/d}{r/d}.\]
For \(r=k\):
\(\sum_{d\mid k}\varphi(d)\binom{k/d}{k/d}=\sum_{d\mid k}\varphi(d)=k\),
so \(N(k,k)=1\) (consistent with \(S_{k,k}=\{\delta_k\}\),
\(B(\delta_k)=\delta_k\)). \(\square\)

\subparagraph{Summary of the answers to cell
1}\label{summary-of-the-answers-to-cell-1}

\begin{itemize}
\tightlist
\item
  \(n=T_k\): every partition reaches \(\delta_k\) (Theorem A), and
  \(\delta_k\) is the unique cyclic partition (Theorem B), forming the
  unique cycle (a fixed point).
\item
  \(n=T_{k-1}+r\), \(1\le r\le k\): the cyclic partitions are exactly
  the \(\binom kr\) partitions \(\lambda^\varepsilon\),
  \(\varepsilon\in\{0,1\}^k\) of weight \(r\) (Theorem B); \(B\) acts on
  them by the rotation
  \(\varepsilon\mapsto(\varepsilon_k,\varepsilon_1,\dots,\varepsilon_{k-1})\)
  (Lemma 1); the number of cycles is
  \(\frac1k\sum_{d\mid\gcd(k,r)}\varphi(d)\binom{k/d}{r/d}\) (Theorem
  C).
\end{itemize}

\emph{Examples.} \(n=8=T_3+2\) (\(k=4,r=2\)): cyclic partitions
\(\lambda^\varepsilon\) for
\(\varepsilon\in\{1100,1010,1001,0110,0101,0011\}\):
\((4,3,1),(4,2,2),(4,2,1,1),(3,3,2),(3,3,1,1),(3,2,2,1)\); cycles:
\(\{1100,0110,0011,1001\}\) and \(\{1010,0101\}\), so
\(N(4,2)=\frac14(\binom42+\varphi(2)\binom21)=2\).

\emph{Independent exact check (not part of the proof).} The script
\texttt{verifica\_cella1.py} (exact integer arithmetic) enumerates all
partitions of every \(n\le40\), computes the cyclic partitions and the
number of cycles by brute force, and confirms they equal \(S_{k,r}\) and
\(N(k,r)\); it also checks Lemma 2 on all compositions of \(n\le14\).
Runtime 7.0 s.


% ===== CAMPO: 3. Verification: instructions, dependencies, timings =====
Python 3 standard library only. Scripts (also saved in
\texttt{runs/p3\_c1/sandbox/} and re-run in
\texttt{runs/p3\_c1/verifica/}): - code\_1 (python, rigor
\texttt{exact}): Exact brute-force check for n=1..40: cyclic partitions
equal S\_\{k,r\}; number of cycles equals the necklace formula; for
n=T\_k the only cyclic partition is delta\_k. Also checks the potential
Lemma 2 (monotonicity and exact equality condition) on all compositions
of n\textless=14. Finite set covered: all partitions of n\textless=40,
all compositions of n\textless=14. Wall-clock 7.0 s. File:
runs/p3\_c1/sandbox/verifica\_cella1.py

Trusted re-runs by the orchestrator (exit code, wall clock, output): -
orchestrator re-ran code\_1.py (python3, clean copy of the researcher
sandbox): exit 0 in 7.0s; stdout: `OK: n=1..40 e potenziale su
composizioni di n\textless=14, tempo 7.0s'; stderr: '\,'


% ===== CAMPO: 4. Sources and contribution =====
\begin{itemize}
\tightlist
\item
  arXiv:1503.00885 (Drensky, survey; abstract only --- context:
  attribution of the characterisation to Brandt 1982)
\item
  arXiv:2607.17194 (Meštrović, survey; abstract only --- context)
\item
  arXiv:1101.1546 (Hart--Khan--Khan, exposition of Toom's proof;
  abstract only --- context)
\item
  arXiv:2208.14496 (Pham; abstract only --- context: necklace
  parametrisation of orbits)
\item
  Position with respect to the literature: The cell's statement is
  classical: 1503.00885 (Drensky) and 2607.17194 (Meštrović) survey it
  and attribute the characterisation of cyclic partitions and the
  necklace count of cycles to Brandt (1982); 1101.1546
  (Hart--Khan--Khan) expounds Toom's proof of convergence to δ\_k for
  triangular n; 2208.14496 (Pham) uses Brandt's necklace parametrisation
  of orbits. I have not read the full papers, only the abstracts;
  nothing from them is used as a hypothesis. My proof follows the
  general strategy of ``diagonal invariant + monotone potential'' but
  the hole-chasing lemma and all details are written out in full here,
  so the result is proved, not cited.
\item
  Contribution: the proof above is written out in full by the team's
  Researcher and checked by two independent judges and by a human.
\end{itemize}


% ===== CAMPO: 5. Limits and unresolved parts =====
Gaps declared by the author (all accepted by the judges): - Lemma 4
(hole lemma) is the most delicate step; the case analysis A/B on the
card's column is written out, but the referee should re-check the index
bookkeeping (in particular that in Case B the card wraps exactly once
before the hole returns to column d-1, and that u' = u-(D-d) is ≥ 1). -
Burnside's lemma and the orbit structure of a rotation on Z\_k are
proved briefly inline (orbit--stabiliser is used without proof); these
are textbook facts.

